% !TeX root = ../../Book3.tex
% 纲要标题：### 20.4 奇点比较
% 纲要位置：计划纲要.md，第 2134 行。

\section{奇点比较}
\label{b3:sec:2004}


% 纲要内容：
%
% * 正规但非正则的例子
% * Cohen–Macaulay 但非正规的例子
% * 正规化与消解奇点不是同一操作
%

\subsection{正规但非正则的例子}
\label{b3:subsec:2004:01}

\begin{example}\label{b3:exm:normal-quadric-cone}
设 \(k\) 为任意域，
\[
A=k[x,y,z]_{(x,y,z)}/(xy-z^2),\qquad\mathfrak m=(\bar x,\bar y,\bar z).
\]
则 \(A\) 是二维正规局部整环，但不是正则局部环。
\end{example}
\begin{proof}
多项式 \(xy-z^2\) 在 \(k[y,z][x]\) 中本原，且在 \(k(y,z)[x]\) 中为一次多项式，故不可约；多项式环为 UFD，因而所得商为整环。它由三维正则局部环商去一个非零因子得到，所以维数二且为 CM 环，特别满足 \(S_2\)。

若 \(\mathfrak p\) 高度一，它不能同时包含 \(\bar x,\bar y\)，否则也包含 \(\bar z\)，便是高度二的 \(\mathfrak m\)。在 \(\bar x\) 可逆处可消去 \(\bar y=\bar z^2/\bar x\)，环成为 \(k[x,x^{-1},z]\) 的局部化；\(\bar y\) 可逆时同样处理。这些环是 UFD，因而高度一局部环是 DVR。高度零处为分式域，所以 \(R_1\) 成立，Serre 判据使 \(A\) 正规。

关系没有一次项，故 \(\mathfrak m/\mathfrak m^2\) 的基为 \(\bar x,\bar y,\bar z\)，维数三，大于 \(\dim A=2\)，所以 \(A\) 不正则。
\end{proof}

这个证明不限制特征；即使特征为二，消去计算与余切空间计算都不改变。正规性允许这个余维二处的奇点存在。

\subsection{Cohen–Macaulay 但非正规的例子}
\label{b3:subsec:2004:02}

一维尖点局部环
\[
R=k[t^2,t^3]_{(t^2,t^3)}\cong k[x,y]_{(x,y)}/(y^2-x^3)
\]
为超曲面环，因而为 CM 环。分式 \(t=y/x\) 满足首一方程 \(T^2-x=0\)，却不在 \(R\) 中：若 \(t=a/b\)，其中 \(a,b\in k[t^2,t^3]\) 且 \(b(0)\ne0\)，则 \(tb\) 的一次项非零，而 \(a\) 的一次项为零，矛盾。因此它不正规。这个计算对任意特征都成立。

\ProseParagraphBreak
交叉直线 \(k[x,y]_{(x,y)}/(xy)\) 同样为 CM 超曲面环，但不是整环，故不是正规局部环。其正规化把两个分支分开；这与尖点中添加同一分式域里的整元素属于不同的几何情形。两者在 Serre 判据中都于高度一的闭点失败。

\ProseParagraphBreak
另一方面，二维正规局部环必为 CM 环：维数为二时 \(S_2\) 恰好要求闭点深度为二。这个结论有明确的维数限制；在更高维中，\(S_2\) 只给深度至少二，不能据此宣称所有正规环都为 CM 环。

\subsection{正规化与奇点消解}
\label{b3:subsec:2004:03}

正规化将整环补成其分式域中的整闭环，或者对约化环的各分支分别进行这一过程。它针对首一方程，而正则性针对局部极大理想的生成结构。\cref{b3:exm:normal-quadric-cone}的环已经正规，其正规化就是自身，原点的非正则性却仍然存在。

\ProseParagraphBreak
在几何中，\term[qidianxiaojie]{奇点消解}[resolution of singularities]寻求由正则空间到原空间的适当双有理映射；“适当”通常要求固有性。这不是仅在同一分式域中添加整元素，也不同于本编自由消解的“消解”。本章只比较二者的目标，不在这里定义完整的概形态射或证明奇点消解存在性。

\ProseParagraphBreak
曲线容易掩盖这个区别：一维 Noether 正规局部整环是 DVR，所以曲线的正规化在局部正则性上已经起了很大作用；二维锥面则直接显示，正规之后仍可能需要改变几何空间。还应区分正则与基域上的光滑，非完美基域上的差别已在\secref{b3:sec:1805}讨论。
